\scanpage{94}
% archive foliation: 96
Pour une courbe elliptique \uline{conforme} donnée, définie
par un réseau \inserted{$\pi$} dans un espace vectoriel \unsure{complexe|conforme}\inserted{$V$} de dim $1$,
les structures réelles qu'on peut mettre dessus correspondent
aux \deleted{\unsure{involutions} $\sigma$}\inserted{droites réelles de $V$}, telles que la
\deleted{\unsure{symétrie}}\inserted{réflexion} orth.
p.r. à $D$ conserve $\pi$. Peut-il y en avoir plusieurs ?
Si oui, les structures réelles correspondantes sont elles isom.
(i.e. conjuguées de l'une d'elles par les autos complexes
de la courbe) ? \deleted{En tous cas, si $\sigma,\sigma'$ sont les deux \unsure{involutions des}}

\deleted{Cas de} Dans tous les cas, soit $G$ le groupe des autos
\uline{conformes} de la courbe elliptique \unsure{conforme}\inserted{$X$ définie par $(\pi,c)$}.
Pour que \deleted{\unsure{l'on}\unsure{ remarque}}\inserted{$X$} admette une structure réelle, il \uline{faut} en
tous cas que $G\overset{\mathrm{det}}\longrightarrow\pm1$ soit surjectif, en autres termes,
on a
\diagram{\begin{tikzcd}[ampersand replacement=\&,column sep=large,row sep=small]
1\arrow[r,dash]\&G^+\arrow[r]\arrow[d,equal]\&G\arrow[r]\&\pm1\arrow[r,dash]\&1\\
\&\operatorname{Aut}(X)\&\&\&
\end{tikzcd}}

Mais un élément $\sigma$ de $G$ de déterminant $-1$
est une transformation orth. de déterminant $-1$,
c'est donc une réflexion, donc c'est une structure
réelle sur la courbe elliptique : les courbes elliptiques
qui admettent une structure réelles sont donc celles
isom. à la conjuguée, i.e. celles dont l'invariant
modulaire est \uline{réel}. L'ens \deleted{de ces structures
(pour $X$ fixée) \inserted{s'identifie à $G^-$ et} est un torseur sous $G^+$, qui est}
\inserted{$G^+$ est cyclique d'ordre} $n=2,4$ ou $6$, et $G$ \deleted{\unsure{se trouve}\unsure{ par}\unsure{ g}\unsure{ 1}\unsure{ 0}}
avec sa structure d'extension,
est isomorphe au groupe diédral $D_n$. L'un des
